\documentclass[11pt]{amsart}
\usepackage[T1]{fontenc}
\usepackage{lmodern}
\usepackage{amsmath,amssymb,amsthm}
\usepackage[margin=1.1in]{geometry}
\usepackage{booktabs}
\usepackage{enumitem}
\usepackage{hyperref}

\newtheorem{theorem}{Theorem}[section]
\newtheorem{lemma}[theorem]{Lemma}
\newtheorem{proposition}[theorem]{Proposition}
\newtheorem{corollary}[theorem]{Corollary}
\newtheorem{conjecture}[theorem]{Conjecture}
\newtheorem{question}[theorem]{Question}
\theoremstyle{definition}
\newtheorem{definition}[theorem]{Definition}
\newtheorem{remark}[theorem]{Remark}
\newtheorem{example}[theorem]{Example}

\newcommand{\cD}{\mathcal D}
\newcommand{\lean}[1]{\quad{\footnotesize[Lean: \texttt{#1}]}}
\newcommand{\PTM}{\mathrm{PTM}}
\newcommand{\STL}{\mathrm{STL}}
\newcommand{\ind}[1]{[\,#1\,]}

\title[Two terminal edges at once]{Two terminal edges at once:\\
the terminal-matching property of Lov\'asz decompositions}
\author{Justin Williams}
\date{November 2026 (working draft)}

\begin{document}
\begin{abstract}
A Lov\'asz decomposition of a graph $G$ whose degrees are all odd is a decomposition into $|G|/2$
paths in which every vertex is the end of exactly one path. For a matching $F$ of $G$ we ask for
a Lov\'asz decomposition in which every edge of $F$ is the terminal edge of its path, at one of
its ends, and no two edges of $F$ lie on the same path. For one edge this is the prescribed end
edge theorem. For two edges it is equivalent to the star-terminal form of the first deletion test
of the spine of Gallai's conjecture at a vertex of degree four, which is how the question arose.

\emph{Statements.} The property fails for two edges in $K_4$ with a perfect matching (this is the
wheel $W_4$ with its diagonal pairing), and for three edges in the prism with its rungs. By
exhaustive computation these are essentially the only obstructions on simple graphs: for two
edges, $K_4$ is the only failure among all instances on at most $8$ vertices (order $10$: see
Section~\ref{sec:comp}); for three edges and for a prescribed end, every failure other than $K_4$
has $G-F$ disconnected. When the edges of $F$ may be parallel to edges of $G$, every failure
other than $K_4$ has a cut vertex in $G-F$.

\emph{Theorems.} We prove a \emph{walk theorem}: attach a leaf in place of the first edge, walk
the resulting even vertex along a path to the second edge, and restore the path edge by edge with
Fan's transfer lemma inside a two-leaf graph, where both prescriptions have become separations of
vertex sets. This reduces two terminal edges to a separation problem with a single even vertex,
which the sparse separation theorem solves when the special vertices are locally sparse. As a
consequence, a graph with all degrees odd has the property for a pair of disjoint edges in one
component, even with one end prescribed, whenever no cycle of length at most four passes through
an end of the two edges; in particular for every pair when the girth is at least five. Lifting
back gives the star-terminal condition, hence the first deletion test, for every vertex $z$ of
degree four with a pairing $ab,cd$ of $N(z)$ such that no cycle of length at most four of the
lifted graph $(J-z)+ab+cd$ passes through a neighbour of $z$. All theorems are formalised in Lean~4 with
Mathlib, without \texttt{sorry} and with only the standard axioms.
\end{abstract}
\maketitle

\section{Introduction}\label{sec:intro}

All graphs are finite and simple unless stated otherwise. A \emph{path decomposition} of $G$ is
a set of paths whose edge sets partition $E(G)$. For a decomposition $\cD$ and a vertex $v$ let
$b_\cD(v)$ be the number of paths of $\cD$ with an end at $v$; always $b_\cD(v)\equiv d(v)\pmod
2$. Following~\cite{RD} we call $\cD$ a \emph{Lov\'asz decomposition} if $b_\cD(v)=d(v)\bmod 2$
for every $v$. Lov\'asz~\cite{L68} proved that every graph with at most one vertex of positive
even degree has one. When all degrees are odd (or zero), every non-isolated vertex is the end of
exactly one path, and for such a vertex $x$ we write $Q_x$ for that path. An edge $xy$ is
\emph{terminal at $x$} if $xy\in E(Q_x)$, that is, if $xy$ is the edge of $Q_x$ at its end $x$.

\begin{theorem}[prescribed end edge~\cite{MS}]\label{thm:endedge}
If $G$ has no vertex of positive even degree, every edge $xy$ is terminal at $x$ in some Lov\'asz
decomposition.
\end{theorem}

This paper is about prescribing several terminal edges at once.

\begin{definition}\label{def:ptm}
Let $G$ have no vertex of positive even degree and let $F$ be a matching of $G$. Then
$\PTM(G,F)$ is the statement that some Lov\'asz decomposition of $G$ has every edge of $F$
terminal at one of its ends, with no two edges of $F$ on the same path. For $e\in F$ and an end
$x$ of $e$, $\PTM(G,F;e@x)$ additionally requires $e$ to be terminal at $x$. \emph{$\PTM_m$} is
the statement ``$\PTM(G,F)$ for every $G$ and every $F$ with $|F|=m$'', and $\PTM_m(r)$ the same
with one end prescribed.
\end{definition}

For $m=1$ this is Theorem~\ref{thm:endedge}. Prescribing all ends (``PEE$_m$'') is false for
$m\ge2$: it fails for $36$ of the $824$ end assignments of order six and $420$ of the $52{,}100$ of
order eight (Section~\ref{sec:comp}). In the data of~\cite[Table~3]{STL} every failure has two
adjacent prescribed ends; if these also have a common neighbour, the transitions close a
triangle. The free choice of ends is therefore essential.

\subsection*{Motivation: the spine}
The spine of the project~\cite{SP} reduces Gallai's conjecture to two deletion tests on
$3$-connected graphs. In the first, $J$ has a unique even vertex $z$, and for every $r\in N(z)$
the graph $J-z$ must have a decomposition into at most $(|J|-1)/2$ paths with two ends at $r$.
The star-terminal programme~\cite{STL} reaches this test from inside Lov\'asz's setting: a
Lov\'asz decomposition of $J$ is \emph{star-terminal} at $z$ if every path through $z$ has the
form $x\,z\,y\cdots$ with $x$ an end, and $\STL(J,z,r)$ asks for one in which $r$ is on a long
side (or some path through $z$ has length two). By~\cite[Theorem~3.2]{STL}, $\STL(J,z,r)$
implies the first test at $r$ (formalised as \texttt{test1\_of\_stl}).

Fix a perfect matching $M$ of $N(z)$ whose pairs are not adjacent in $J$, and let
$G=(J-z)+M$. Replacing each edge $xy\in M$ by the path $x\,z\,y$ turns a Lov\'asz decomposition
of $G$ in which every edge of $M$ is terminal and no path has two of them into a star-terminal
Lov\'asz decomposition of $J$ with transitions $M$ at $z$, and conversely~\cite[Proposition~4.4]{STL}.
An edge $xy\in M$ terminal at $x$ lifts to $x\,z\,y\cdots$, so $x$ is short and $y$ is long. Hence
\[
\PTM(G,M;e@s)\ \Longrightarrow\ \STL(J,z,r)\qquad(e=rs\in M).
\]
At $d(z)=4$, two edges are prescribed. This is the composition problem of~\cite[\S X.4a]{CONN}:
the single-transition tools (separation, the leaf trick, the prescribed end edge, the end-edge
theorem EE at the even vertex~\cite{PE}) each create one even vertex per prescription, and two
even vertices are beyond them.

\subsection*{Results}
\begin{enumerate}[label=(\arabic*)]
\item \emph{The right statements} (Section~\ref{sec:statements}). $K_4$ with a perfect matching
violates $\PTM_2$ (Example~\ref{ex:K4}); it is the wheel $W_4$ with its diagonal pairing, the only
known failure of the pairing-prescribed form of STL with $J$ $3$-connected. The prism with its
rungs violates $\PTM_3$. Exhaustive computation (Section~\ref{sec:comp}) supports
Conjecture~\ref{conj:ptm}: $\PTM_2$ holds for every simple $G$ except $K_4$, and $\PTM_m(r)$ holds
whenever $G-F$ is connected and $G\ne K_4$. When edges of $F$ are allowed to be parallel to edges
of $G$ (pairs of $M$ adjacent in $J$), every failure found other than $K_4$ has a cut vertex in
$G-F$ (Conjecture~\ref{conj:multi}).
\item \emph{The walk theorem} (Theorem~\ref{thm:walk}). Two terminal edges reduce to a
separation problem in a graph with one even vertex: three pairs of vertices must be separated in
a Lov\'asz decomposition of a base graph obtained by deleting a path.
\item \emph{The sparse form} (Theorem~\ref{thm:sparse}). That separation problem is solved by the
sparse separation theorem of~\cite{MS} when four special vertices are sparse, up to one
exceptional vertex, which is removed by one more Fan transfer (Lemma~\ref{lem:fix}).
\item \emph{Local girth five} (Theorem~\ref{thm:girth}). If $G$ has all degrees odd, $F$ is a
pair of disjoint edges in one component, and no triangle and no $4$-cycle passes through an end
of an edge of $F$, then $\PTM(G,F;e@x)$ holds for every choice of $e\in F$ and of its end $x$. In
particular this holds for every $F$ when $G$ is connected and has girth at least five.
\item \emph{STL and the first test} (Theorem~\ref{thm:stl}, Corollary~\ref{cor:test1}). If $z$ has
degree four, $N(z)=\{a,b,c,d\}$ with $ab,cd\notin E(J)$, and no triangle and no $4$-cycle of
$(J-z)+ab+cd$ passes through $a$, $b$, $c$ or $d$ (and $c$ is reachable from $a$ in it), then
$\STL(J,z,r)$ holds for every $r\in N(z)$, and the first deletion test holds at $z$.
\item \emph{Lean} (Section~\ref{sec:lean}). Items (2)--(5) are formalised in
\texttt{lean/GallaiLean/TerminalMatching.lean}, sorry-free with only the standard axioms, and
connected to \texttt{test1\_of\_stl}.
\end{enumerate}

\subsection*{Honest status}
$\PTM_2$ in general is not proved. The walk theorem is a reduction, and its base hypothesis holds
in $93.6\%$ of all instances on eight vertices when only shortest walks are used
(Section~\ref{sec:comp}); the sparse form that we can prove covers local girth five, and
$17.6\%$ of the (mostly dense) eight-vertex instances. The consequence for the spine needs sparse
graphs: for random $3$-connected graphs with degrees $4,3,\dots,3$ the local hypothesis of
Corollary~\ref{cor:test1} holds for some pairing in $2\%$ of the graphs of order $11$ and $70\%$
of those of order $51$ (global girth five: $2\%$ and $5\%$; Section~\ref{sec:comp}). Its value
is that it is the first unconditional two-prescription result in this programme, and that the
mechanism (walk the even vertex, restore inside the leaf graph) is general.

\section{Statements, examples and conjectures}\label{sec:statements}

\begin{example}[$K_4$]\label{ex:K4}
Let $G=K_4$ on $\{a,b,c,d\}$ and $F=\{ac,bd\}$. A Lov\'asz decomposition has two paths with three
edges each, so both are Hamiltonian paths. A Hamiltonian path in which $ac$ is terminal at $a$
reads $a\,c\,u\,v$ with $\{u,v\}=\{b,d\}$, so its last edge is $bd$; the same holds with the roles
of $a$ and $c$ exchanged. So $ac$ and $bd$ lie on one path, and $\PTM(K_4,F)$ fails. Here
$J=(G-F)+z(V(F))$ is the wheel $W_4$ and $F$ its diagonal pairing, the exception
in~\cite[Example~4.5]{STL}. The pairs of $F$ are not adjacent in $J$; this is a simple-graph
failure.
\end{example}

\begin{example}[the prism]\label{ex:prism}
Let $G$ be the prism, two triangles $x_1x_2x_3$ and $y_1y_2y_3$ joined by the rungs
$F=\{x_iy_i\}$. Then $G-F$ is two triangles, which need four paths, while a decomposition as in
$\PTM(G,F)$ lifts, by deleting $z$ from the corresponding star-terminal decomposition
of $J=(G-F)+z(V(F))$, to a decomposition of $G-F$ into at most $(7-1)/2=3$
paths~\cite[Lemma~3.1]{STL}. So $\PTM(G,F)$ fails, with $|F|=3$.
\end{example}

The second example generalises: whenever $\PTM(G,F)$ holds, $G-F$ has a decomposition into at
most $|G|/2$ paths (delete $z$ from the lifted decomposition). So a disconnected $G-F$ with too
many components is an obstruction, and connectivity of $G-F$ (of $J-z$ in~\cite{STL}) is the
natural hypothesis.

\emph{Parallel edges.} For a pairing $M$ of $N(z)$ with a pair $xy$ adjacent in $J$, the graph
$(J-z)+M$ has two parallel edges $xy$. We call this the \emph{multigraph case}: $G'=G-F$ is
simple, $F$ is a perfect matching of the even-degree vertices of $G'$, and an edge of $F$ may
join two adjacent vertices of $G'$. The oracle of Section~\ref{sec:comp} works with $J$
directly, so it covers both cases.

\begin{conjecture}[$\PTM$, simple case]\label{conj:ptm}
Let $G$ be a simple graph with all degrees odd and $F$ a matching of $G$ such that $G-F$ is
connected. Then $\PTM(G,F;e@x)$ holds for every $e\in F$ and every end $x$ of $e$, unless $G=K_4$
and $|F|=2$. Moreover $\PTM(G,F)$ holds for $|F|=2$ even when $G-F$ is disconnected, unless the
component of $G$ containing $F$ is $K_4$.
\end{conjecture}

\begin{conjecture}[$\PTM$, multigraph case]\label{conj:multi}
In the multigraph case, if $G-F$ is $2$-connected and $G\ne K_4$ (as in Example~\ref{ex:K4}),
then $\PTM(G,F;e@x)$ holds for every $e$ and $x$.
\end{conjecture}

If $J$ is $3$-connected then $J-z$ is $2$-connected, so Conjecture~\ref{conj:multi} implies the
pairing-prescribed STL for $3$-connected $J\ne W_4$~\cite[Conjecture~7.4]{STL}, and with it STL
and the first deletion test.

\section{Tools}\label{sec:tools}

For a decomposition $\cD$ and a vertex $g$ let $Z_\cD(g)=\{w\in N(g):b_\cD(w)=0\}$.

\begin{lemma}[Fan's lemma with vertex sets~\cite{Fan05},~\cite[Lemma~2.2]{RD}]\label{lem:fan}
Let $g\ne t$ be adjacent in $K$, and let $\cD$ decompose $K-gt$ with $b_\cD(t)>|Z_\cD(g)|$, where
$Z$ is taken in $K-gt$. Then $K$ has a decomposition $\cD^+$ with
$b_{\cD^+}=b_\cD+\mathbf 1_g-\mathbf 1_t$, every path of which has its vertex set inside
$V(Q)\cup\{g\}$ for some path $Q$ of $\cD$.\lean{fan\_single\_supp}
\end{lemma}

\begin{theorem}[sparse separation~\cite{MS}]\label{thm:sparsesep}
Let $K$ have at most one vertex of positive even degree, let $X\subseteq V(K)$ be \emph{sparse}
(every vertex of $K$ has at most one neighbour in $X$), and let $\mathcal A$ be a family of
subsets of $X$, each of which is packed (its vertices are pairwise non-adjacent without common
neighbours). Then some Lov\'asz decomposition of $K$ meets every member of $\mathcal A$ in at most
one vertex per path.\lean{separation\_sparse}
\end{theorem}

\emph{Leaves.} For a graph $K$ and a vertex $r$, $K+\ell r$ is $K$ with a new vertex $\ell$
joined to $r$. For an edge $xy$ of $K$, the \emph{leaf graph} $K\langle x,y\rangle$ is
$(K-xy)+\ell y$; the leaf $\ell$ stands for $x$.

\begin{lemma}[lift and merge~\cite{PE,MS}]\label{lem:leaf}
\begin{enumerate}[label=\textup{(\alph*)}]
\item If $P\in\cD$ ends at $r$, extending $P$ by $r\ell$ gives a decomposition of $K+\ell r$ with
one end fewer at $r$, one end at $\ell$, the other end counts of $\cD$, and every path either inside
$V(P)\cup\{\ell\}$ or with the vertex set of a path of $\cD$.\lean{PathDecomp.liftLeaf}
\item Let $\cD_H$ decompose $K\langle x,y\rangle$ with no path containing both $x$ and $\ell$.
Identifying $\ell$ with $x$ gives a decomposition $\cD$ of $K$ with
$b_\cD(x)=b_{\cD_H}(x)+b_{\cD_H}(\ell)$ and the other end counts unchanged; each path of $\cD$ is the
image of a path of $\cD_H$, with the same ends and edges up to $\ell\mapsto x$.\lean{leaf\_merge}
\end{enumerate}
\end{lemma}

In particular, in (b) the path through $y\ell$, which ends at the leaf $\ell$, becomes a path
ending at $x$ whose terminal edge is $yx$. In the leaf graph, ``$xy$ is terminal at $x$'' is
the vertex-set property ``$x$ and $\ell$ are separated'', and Lemma~\ref{lem:fan} preserves it
whenever its target is neither $x$ nor $\ell$. This is the observation behind the protected-end
lemma of~\cite[Lemma~4.1]{PE}, and the engine of everything below.

\begin{lemma}[walking back]\label{lem:walkback}
Let $w_0w_1\cdots w_k$ $(k\ge0)$ be a path in $K$, let $K_0=K-\{w_0w_1,\dots,w_{k-1}w_k\}$, and
let $\cD$ decompose $K_0$. Suppose that every vertex $v\ne w_0$ adjacent in $K$ to one of
$w_0,\dots,w_{k-1}$ has $b_\cD(v)\ge1$. Then $K$ has a decomposition $\cD'$ with
$b_{\cD'}=b_\cD+\mathbf 1_{w_0}-\mathbf 1_{w_k}$, every path of which lies in a path of $\cD$
together with some of $w_0,\dots,w_{k-1}$.\lean{walk\_back}
\end{lemma}

\begin{proof}
Induction on $k$; for $k=0$ take $\cD'=\cD$. For $k\ge1$ let
$K_1=K-\{w_1w_2,\dots,w_{k-1}w_k\}$, so $K_0=K_1-w_0w_1$. Every neighbour of $w_0$ in $K_0$ has an
end, so $Z_\cD(w_0)=\emptyset$, and $b_\cD(w_1)\ge1$ because $w_1$ is adjacent to $w_0$.
Lemma~\ref{lem:fan} with target $w_0$ and $t=w_1$ gives $\cD_1$ on $K_1$ with
$b_{\cD_1}=b_\cD+\mathbf 1_{w_0}-\mathbf 1_{w_1}$, whose paths lie in paths of $\cD$ plus $w_0$.
The hypothesis holds for $\cD_1$ and the path $w_1\cdots w_k$: a vertex $v\ne w_1$ adjacent to one
of $w_1,\dots,w_{k-1}$ either is $w_0$, which now has an end, or keeps its end count. Induction
gives $\cD'$ on $K$, and the end counts and vertex sets compose.
\end{proof}

\section{The walk theorem}\label{sec:walk}

Throughout this section $G$ has no vertex of positive even degree, and $ab$, $xy$ are disjoint
edges of $G$. We aim at $ab$ terminal at $a$ and $xy$ terminal at $x$, on different paths.

\begin{definition}\label{def:walk}
An \emph{admissible walk} is a path $P=p_0p_1\cdots p_k$ of $G$ with $k\ge1$, $p_0=y$, $p_k=a$,
$b,x\notin V(P)$, and $p_i\notin N(x)$ for $0<i<k$. Its \emph{base graph} is
\[
M=\bigl((G-ab)+\ell_1b\bigr)-xy-E(P),
\]
where $\ell_1$ is a new leaf at $b$ (it stands for $a$).
\end{definition}

The path $P$ uses neither $ab$ nor $xy$, because $b,x\notin V(P)$. In $M$ the degrees of $a$ and
$y$ drop by two, those of the inner vertices of $P$ drop by two, that of $x$ drops by one, and
$\ell_1$ has degree one. So $x$ is the only vertex of $M$ that can have positive even degree, and
$d_M(v)\equiv d_G(v)+\ind{v=x}\pmod 2$.\lean{baseG\_parity}

\begin{theorem}[walk theorem]\label{thm:walk}
Let $P$ be an admissible walk with base graph $M$. If $M$ has a Lov\'asz decomposition in which no
path contains both vertices of any of the pairs $\{x,y\}$, $\{a,\ell_1\}$, $\{\ell_1,y\}$, then $G$
has a Lov\'asz decomposition in which $ab$ is terminal at $a$, $xy$ is terminal at $x$, and the two
paths are different.\lean{ptm2\_of\_base}
\end{theorem}

\begin{proof}
Let $\cD_0$ be the given decomposition of $M$. Every non-isolated vertex of $M$ other than $x$ has
odd degree, hence one end, and $b_{\cD_0}(x)=0$.

\emph{Step 1: the second leaf.} Let $Q_y$ be the path of $\cD_0$ ending at $y$, and extend it by a
new leaf $\ell_2$ at $y$ (Lemma~\ref{lem:leaf}(a)). The result $\cD_1$ decomposes $M+\ell_2y$,
with $b_{\cD_1}(y)=0$, $b_{\cD_1}(\ell_2)=1$ and the other end counts of $\cD_0$. Every path of
$\cD_1$ is either $Q_y\ell_2$ or a path of $\cD_0$. Since $Q_y$ contains $y$, the hypothesis gives:
no path of $\cD_1$ contains $\{x,\ell_2\}$, $\{\ell_1,\ell_2\}$ or $\{a,\ell_1\}$.

\emph{Step 2: walking back.} Let $H$ be the two-leaf graph $\bigl((G-ab)+\ell_1b-xy\bigr)+\ell_2y$.
Then $M+\ell_2y=H-E(P)$. Apply Lemma~\ref{lem:walkback} in $H$ to the path $y=p_0,\dots,p_k=a$. A
vertex with no end in $\cD_1$ is $x$, $y=p_0$, or isolated in $G$ (every other vertex of $G$ has
odd degree in $M$, by the parity computation above, and the leaves have one end each). A vertex
isolated in $G$ is adjacent to no $p_i$. No vertex $p_i$ with $0\le i<k$ is
adjacent to $x$ in $H$: for $i=0$ the edge $xy$ has been removed, and for $0<i<k$ this is
admissibility. So the hypothesis of Lemma~\ref{lem:walkback} holds. It gives a decomposition
$\cD_2$ of $H$ with $b_{\cD_2}(a)=0$, $b_{\cD_2}(y)=1$, $b_{\cD_2}(x)=0$, one end at each leaf, and
one end at every other non-isolated vertex. Every path of $\cD_2$ lies in a path of $\cD_1$
together with some of $p_0,\dots,p_{k-1}$. These targets are not $x$ (as $x\notin V(P)$), not $a$
($a=p_k$ appears once on $P$), and not a leaf. So the three separations survive:
no path of $\cD_2$ contains $\{x,\ell_2\}$, $\{\ell_1,\ell_2\}$ or $\{a,\ell_1\}$.

\emph{Step 3: merging.} The graph $H$ is the leaf graph $K\langle x,y\rangle$ of
$K=(G-ab)+\ell_1b$, with $\ell_2$ standing for $x$. Lemma~\ref{lem:leaf}(b) merges $\ell_2$ into
$x$: we get a decomposition $\cD_3$ of $K$ with $b_{\cD_3}(x)=0+1$. The path through $y\ell_2$
becomes a path $R_2$ ending at $x$ through $yx$. A path of $\cD_3$ containing $a$ and $\ell_1$
would come from a path of $\cD_2$ containing $a$ and $\ell_1$, since the merge only renames
$\ell_2$ as $x\ne a$; so $a$ and $\ell_1$ are separated in $\cD_3$. Now $K$ is the leaf graph
$G\langle a,b\rangle$, and Lemma~\ref{lem:leaf}(b) merges $\ell_1$ into $a$. It gives a Lov\'asz
decomposition $\cD$ of $G$ in which the image $R_1$ of the path through $b\ell_1$ ends at $a$
through $ba$, and the image of $R_2$ ends at $x$ through $yx$. The end counts are those of a
Lov\'asz decomposition of $G$: $a$ has $0+1$ ends, $x$ has one, and every other vertex keeps its
count.

\emph{Distinctness.} If the two final paths coincided, then $R_2$, as a path of $\cD_3$, would
contain the edge $b\ell_1$, because this is the only preimage of $ba$ in $K$. Then its preimage in
$\cD_2$, the path ending at $\ell_2$, would contain $\ell_1$. That contradicts the separation of
$\ell_1$ and $\ell_2$.
\end{proof}

\begin{remark}
During Step~2 the intermediate graphs have two vertices of even degree without ends: the short
end $x$, and the current target $p_i$ (at the end, the short end $a$). So they are outside
Lov\'asz's setting, and nothing could be obtained there by induction. Fan's lemma does not need
parity, only the end count of its source and the absence of endless neighbours of its target.
This is why the walk can be carried out with both prescriptions frozen as separations. The
admissibility condition $p_i\notin N(x)$ is exactly what keeps the endless short end $x$ out of
$Z(p_i)$.
\end{remark}

\section{The base case: sparse separation with one exception}\label{sec:sparse}

\begin{lemma}[one exceptional vertex]\label{lem:fix}
Let $g$ be the only vertex of $K$ that can have positive even degree, with $d(g)$ even, and let
$v\in N(g)$. If $X$ and $\mathcal A$ satisfy the hypothesis of Theorem~\ref{thm:sparsesep} in
$K-gv$ and $v\notin X$, then $K$ has a Lov\'asz decomposition meeting every member of $\mathcal
A$ at most once per path.\lean{sep\_sparse\_fix}
\end{lemma}

\begin{proof}
In $K-gv$ the vertex $g$ is odd and $v$ is the only vertex that can have positive even degree.
Theorem~\ref{thm:sparsesep} gives a Lov\'asz decomposition of $K-gv$ with the separations; there
$b(g)=1$, $b(v)=0$, and every neighbour of $v$ has odd degree, hence an end. Lemma~\ref{lem:fan}
with target $v$ and $t=g$ restores $gv$. Every new path lies in an old path plus $v\notin X$, so
the separations survive, and the end counts become $b(v)=1$, $b(g)=0$: a Lov\'asz decomposition of
$K$.
\end{proof}

This is the mechanism of~\cite[Theorem~5.1]{PE} (one common neighbour of $u$ and $z$), isolated
as a lemma.

\begin{theorem}[walk theorem, sparse form]\label{thm:sparse}
Let $P$ be an admissible walk with base graph $M$, and let $X=\{x,y,a,\ell_1\}$. Suppose that the
pairs $\{x,y\}$, $\{a,\ell_1\}$, $\{\ell_1,y\}$ are packed and $X$ is sparse, either in $M$ or in
$M-xv$ for one neighbour $v\notin X$ of $x$. Then $G$ has a Lov\'asz decomposition with $ab$
terminal at $a$ and $xy$ terminal at $x$, on different paths.\lean{ptm2\_of\_sparse}
\end{theorem}

\begin{proof}
$x$ is the only possible vertex of positive even degree of $M$, and $d_M(x)$ is even. Apply
Theorem~\ref{thm:sparsesep} or Lemma~\ref{lem:fix}, then Theorem~\ref{thm:walk}.
\end{proof}

\section{Girth five}\label{sec:girth}

\begin{theorem}[local girth five]\label{thm:girth}
Let $G$ have no vertex of positive even degree. Let $ab$ and $cd$ be disjoint edges, with $c$
reachable from $a$, such that no cycle of length three or four passes through $a$, $b$, $c$ or
$d$. Then $G$ has a Lov\'asz decomposition in which $ab$ is terminal at $a$, $cd$ is terminal at
$c$ or at $d$, and $ab$, $cd$ lie on different paths.\lean{ptm2\_localgirth, ptm2\_girth5}
\end{theorem}

The end of $ab$ is prescribed, so for every $r\in\{a,b,c,d\}$ there is such a decomposition in
which $r$ is the long end of its edge: apply the theorem to $(b,a)$ or to $(a,b)$, with the roles
of the two edges exchanged if $r\in\{c,d\}$. In the language of Definition~\ref{def:ptm}:
\emph{$\PTM_2$ and $\PTM_2(r)$ hold for all graphs of girth at least five with all degrees odd},
when the two edges lie in one component; otherwise apply Theorem~\ref{thm:endedge} in each
component. More generally $\PTM(G,F;e@x)$ holds for every pair $F$ none of whose ends lies on a
cycle of length three or four.

\begin{proof}
Choose $u\in\{a,b\}$, $w\in\{c,d\}$ and a walk from $w$ to $u$ of least possible length $k$ over
all four choices. Bypassing cycles gives a path of length at most $k$, hence exactly $k$; call it
$P=p_0\cdots p_k$, $p_0=w$, $p_k=u$. Then $k\ge1$ since the edges are disjoint.

\emph{Orientation.} If $u=a$, we use the walk theorem with short ends $a$ and $x$, where $y=w$ and
$x$ is the other end of $cd$. If $u=b$, the prescribed long end $b$ is the end of $P$; we
exchange the roles of the two edges: the walk theorem is applied to the edge $cd$ with short end
$w$, to the edge $ab$ with $y=b$ and short end $x=a$, and to the reversed path. In both cases we
are in the following situation, which we now treat: $ab$ and $xy$ are disjoint edges, $P$ is a
path from $y$ to $a$ of length $k$, and every walk between $\{a,b\}$ and $\{x,y\}$ has length at
least $k$. We show that Theorem~\ref{thm:sparse} applies; it yields $ab$ terminal at $a$ and $xy$
terminal at $x$. Every triangle and every $4$-cycle excluded below passes through one of $a$, $b$,
$x$, $y$, which is all the hypothesis provides.

\emph{Admissibility.} If $b\in V(P)$, the part of $P$ from $y$ to $b$ is shorter than $k$, since
$b\ne a$. If $x\in V(P)$, the part from $x$ to $a$ is shorter. If an inner vertex $p_i$ were
adjacent to $x$, the walk $x\,p_i\,p_{i+1}\cdots p_k$ would have length $k-i+1$. That is less than
$k$ unless $i=1$, and then $p_1$ would be adjacent to both $x$ and $y=p_0$: a triangle.

\emph{Short distances.} If $b$ were adjacent to $x$ or $y$, some walk between the two edges would
have length one. Then $k=1$, $P=ya$, and $a\,y\,x\,b$ would be a $4$-cycle or $a\,b\,y$ a
triangle. If $ya\in E(G)$ then $k=1$ and $ya$ is the edge of $P$.

\emph{Sparseness.} The only neighbour of $\ell_1$ in $M$ is $b$, and $b$ is adjacent in $M$ to
none of $a$, $x$, $y$. So $b$ has exactly one neighbour in $X$. Let $v\ne b$ be adjacent in $M$ to
two vertices of $X\setminus\{\ell_1\}=\{x,y,a\}$.
\begin{itemize}
\item $v\sim x,y$: a triangle.
\item $v\sim y,a$ by edges of $M$ (not of $P$): the walk $y\,v\,a$ gives $k\le2$. If $k=1$ then
$v\,y\,a$ is a triangle. If $k=2$, then $P=y\,p_1\,a$ with $v\ne p_1$, since the edges $yp_1$ and
$p_1a$ are not in $M$, and $y\,v\,a\,p_1$ is a $4$-cycle.
\item $v\sim x,a$: the walk $x\,v\,a$ gives $k\le2$. If $k=1$ then $a\,v\,x\,y$ is a $4$-cycle;
here $v\ne y$ because $ya\in E(P)$ is not an edge of $M$.
If $k=2$ this can happen (a $5$-cycle $a\,v\,x\,y\,p_1$). Such a $v$ is unique: two of them, $v$
and $v'$, would give the $4$-cycle $x\,v\,a\,v'$. And $v\notin X$: $v\ne y$, because $ya\in E(G)$
forces $k=1$ with $ya\in E(P)$.
\end{itemize}
The vertices of $X$ themselves have no neighbour in $X$ in $M$: $xy$ and $ab$ are deleted, $ya$ is
an edge of $P$ if present, $xa\in E(G)$ would give $k=1$ and the triangle $a\,x\,y$, and
$\ell_1$ is adjacent only to $b$. Packedness: $x$ and $y$ are non-adjacent in $M$ with no common
neighbour (a triangle); $a$ and $\ell_1$ have no common neighbour because $b\not\sim a$ in $M$;
$\ell_1$ and $y$ are non-adjacent, and $b\not\sim y$. If there is no exceptional vertex $v$, $X$
is sparse in $M$. Otherwise $X$ is sparse in $M-xv$, where packedness persists because deleting an
edge only helps. Theorem~\ref{thm:sparse} applies.
\end{proof}

\section{STL and the first deletion test at a vertex of degree four}\label{sec:stl}

Let $J$ have $z$ as its only even vertex, with $N(z)=\{a,b,c,d\}$ and $ab,cd\notin E(J)$. Let
\[
J^\ast=(J-z)+ab+cd,
\]
regarded as a graph on $V(J)$ in which $z$ is isolated. It is obtained by lifting the transitions
$a\,z\,b$ and $c\,z\,d$ at $z$, one after the other.\lean{liftJ, swapG}

\begin{theorem}\label{thm:stl}
If $J^\ast$ has a Lov\'asz decomposition with $ab$ terminal at $a$ and $cd$ terminal at $c$, on
different paths, then $J$ has a star-terminal Lov\'asz decomposition with transitions $a\,z\,b$
and $c\,z\,d$ in which $b$ and $d$ are on long sides, or a path through $z$ has length two. In
particular $\STL(J,z,b)$ and $\STL(J,z,d)$ hold.\lean{stl\_of\_ptm2}
\end{theorem}

\begin{proof}
Every vertex other than $z$ has the same degree in $J^\ast$ as in $J$, and $z$ has degree $0$, so
the decomposition is a Lov\'asz decomposition of $J$-parities. Un-lift $cd$: the path ending at $c$
through $cd$ reads $c\,d\cdots$ and avoids $z$, and replacing $cd$ by $c\,z\,d$ gives a path of
$J^\ast-cd+cz+zd$. All other paths are kept. Then un-lift $ab$ in the same way. Every vertex keeps
its ends. The only paths through $z$ are $a\,z\,b\cdots$ and $c\,z\,d\cdots$, so the
decomposition is star-terminal. The side at $b$ is long unless the path is $a\,z\,b$, which has
length two; likewise at $d$.
\end{proof}

\begin{corollary}\label{cor:test1}
Let $z$ be the only even vertex of $J$, with $N(z)=\{a,b,c,d\}$, $ab,cd\notin E(J)$, and every
neighbour of $z$ of degree at least three. If no triangle and no $4$-cycle of $J^\ast$ passes
through $a$, $b$, $c$ or $d$, and $c$ is reachable from $a$ in $J^\ast$, then $\STL(J,z,r)$ holds for every $r\in N(z)$, and the first
deletion test holds at $z$: for every $r\in N(z)$, $J-z$ has a decomposition into at most
$(|J|-1)/2$ paths with two ends at $r$.\lean{stl\_of\_localgirth, test1\_of\_localgirth}
\end{corollary}

\begin{proof}
$J^\ast$ has no vertex of positive even degree. For $r\in N(z)$, Theorem~\ref{thm:girth} with
the suitable labelling (Section~\ref{sec:girth}) makes $r$ the long end of its pair, and
Theorem~\ref{thm:stl} gives $\STL(J,z,r)$. The first test follows from STL~\cite[Theorem~3.2]{STL}
(\texttt{test1\_of\_stl}).
\end{proof}

In terms of $J$: no cycle of length at most four of $J^\ast$ passes through $a,b,c,d$ exactly
when no such cycle of $J-z$ does, $d_{J-z}(a,b),d_{J-z}(c,d)\ge4$, and
$d_{J-z}(a,c)+d_{J-z}(b,d)\ge3$ and $d_{J-z}(a,d)+d_{J-z}(b,c)\ge3$. (With ``$J-z$ has girth at
least five'' in place of the first condition, this characterises girth five of $J^\ast$.) Each of
the other two pairings gives its own condition, and one of them suffices.

\section{Computations}\label{sec:comp}

\emph{This section is computational; no proof above depends on it.} Scripts are in
\texttt{scripts/}, outputs in \texttt{scripts/out/}; the experiment log is E24 in
\texttt{research/experiments.md}.

\subsection*{Oracle}
\texttt{ptm.c} decides $\PTM(G,F)$, $\PTM(G,F;e@x)$ and the prescribed-end form through
$J=(G-F)+z(V(F))$. It searches for a star-terminal Lov\'asz decomposition of $J$ with allowed
transitions $F$ at $z$ and, for $e@x$, the required terminal edge $xz$. The search engine is a
verbatim copy of \texttt{tools/stl.c}~\cite{STL}, validated there against two brute-force
enumerators: a depth-first search over paths with memoised failures, exact unless a node budget
is given. The only change is that the memo table is reset by generation stamps. Every positive
answer is re-checked here by an independent verifier (paths, edge partition, end counts, star
condition, transitions, prescribed ends); there were no bad witnesses in any run.
Because the oracle works with $J$, the multigraph case (an $F$-edge parallel to an edge of $G$) is
handled without change. For order $10$, the instances $(G,F)$ are reduced up to isomorphism.
\texttt{encode.c} subdivides each $F$-edge by a vertex carrying two pendant leaves, the only
vertices of even degree; nauty's \texttt{labelg} canonicalises the result; and \texttt{decode.py}
recovers $(G-F,F)$.

\subsection*{Results}
\begin{table}[!ht]
\centering\footnotesize
\begin{tabular}{@{}lrrrr@{}}
\toprule
class & instances & PTM fails & PTM$(e@x)$ tests & fails\\
\midrule
simple, $m=2$, $n\le8$ & 13{,}234 & 3 ($K_4$) & 52{,}924 & 0\\
simple, $m=3$, $n\le8$ & 14{,}124 & 2 & 84{,}732 & 15 \\
simple, $m=4$, $n=8$ & 1{,}955 & 0 & 15{,}640 & 14 \\
simple, $m=2$, $n=10$ (up to isomorphism) & 3{,}474{,}599 & \textit{pending} & & \\
cubic, $m=2$, $n=12$ (up to isomorphism) & 3{,}745 & 0 & 14{,}980 & 0\\
cubic, $m=2$, $n=14$ (up to isomorphism) & 46{,}407 & 0 & 185{,}628 & 0\\
girth $\ge5$, $m=2$, $n\le12$ & 1{,}497 & 0 & 5{,}988 & 0\\
multigraph, $G-F$ connected, $n\le8$ & 76{,}632 & 29 & 453{,}964 & 1{,}942\\
\bottomrule
\end{tabular}
\caption{Exhaustive tests. ``Simple'': all connected all-odd $G$ and all matchings $F$ of $m$
edges (all instances, not up to isomorphism, except at $n=10$). PTM$(e@x)$ is tested only on
the instances where PTM holds, so $K_4$ is not counted in that column. Every PTM failure with
$m=3$ and every PTM$(e@x)$ failure in the simple rows has $G-F$ disconnected; the $m=3$ failures are the
prism (order $6$) and one graph of order $8$. ``Multigraph'': all connected $G'$ of order $\le8$
and all perfect matchings $F$ of the even vertices of $G'$ with $G'+F$ connected; every failure
except $K_4$ has a cut vertex in $G'=G-F$ (and $J$ of connectivity at most $2$).}
\label{tab:exh}
\end{table}

\subsection*{Random tests}
Simple instances of orders $12$, $14$ and $16$ ($G(n,p)$ with $p\in[0.15,0.5]$ and parities
repaired; $m=2,3,4$; $F$ random, or each new edge of $F$ adjacent to an earlier one, or at
distance exactly two from it so that the edges share neighbourhoods; $150$ instances per cell):
$4{,}050$ instances and $23{,}950$ end-prescribed queries, no failure. With a node budget,
$56$ PTM and $75$ PTM$(e@x)$ queries were undecided; a re-run of the $120$ instances concerned
with budget $5\cdot10^6$ decided all but $8$ and $23$ of them, every one positively (an
independent brute force of audit~A22 decides $82$ of the $120$, all positively). Multigraph
instances with $G-F$ $2$-connected (Conjecture~\ref{conj:multi}): $120$ random graphs $G'$ for
each order $10$, $12$, $14$ and each $m=2,3$, with $2m$ even vertices and every pairing $F$ of
them; $6{,}480$ instances, $4{,}978$ with an $F$-edge parallel to an edge of $G'$, no failure
($48$ PTM and $86$ PTM$(e@x)$ queries undecided within the budget).

\subsection*{How often Corollary~\ref{cor:test1} applies}
\texttt{localgirth\_stat.py}: random $3$-connected graphs $J$ with degree sequence $4,3,\dots,3$
(configuration model, simple graphs kept), $300$ per order; a graph counts if some pairing of
$N(z)$ into non-adjacent pairs satisfies the hypothesis.
\begin{center}\footnotesize
\begin{tabular}{@{}lrrrrr@{}}
\toprule
order & $11$ & $15$ & $21$ & $31$ & $51$\\
\midrule
local hypothesis (Corollary~\ref{cor:test1}) & $5$ & $29$ & $73$ & $138$ & $210$\\
$J^\ast$ of girth $\ge5$ & $5$ & $3$ & $6$ & $11$ & $14$\\
\bottomrule
\end{tabular}
\end{center}

\subsection*{How far the walk theorem reaches}
\texttt{walk.c} tests, for every instance with $|F|=2$, all eight choices of which edge is $ab$
and which ends are short, and all admissible walks of least length (or up to a given excess). It
decides whether the base graph has a Lov\'asz decomposition with the three separations (an exact
search, the oracle with a separation constraint), or, in the second mode, whether the sparse
hypotheses of Theorem~\ref{thm:sparse} hold.

\begin{center}\footnotesize
\begin{tabular}{@{}lrrrr@{}}
\toprule
instances ($n\le8$, all $2$-matchings) & total & no admissible walk & base hypothesis & sparse form\\
\midrule
$n=6$ & 206 & 1 & 160 & 31\\
$n=8$ & 13{,}025 & 68 & 12{,}193 & 2{,}288 \\
$n=8$, walks up to $3$ longer than shortest & 13{,}025 & 68 & & 3{,}093\\
\bottomrule
\end{tabular}
\end{center}

So the reduction to a separation problem is strong, much stronger than the part of it that is
proved. But it is lossy: in $832$ instances of order eight the base hypothesis fails for every
orientation and every shortest admissible walk ($764$ of them have such a walk, $68$ have none),
although $\PTM_2$ holds in all of them; allowing walks up to two longer than shortest covers no
further instance. So the base separation problem (three pairs forming a path $x-y-\ell_1-a$ of
constraints, one even vertex) is not always solvable. When the three pairs are packed in $M$, it
is a case of the class-separation questions of~\cite{MS} for the union of two matchings, where no
counterexample is known; all $11{,}810$ base graphs of order eight with packed pairs are
separable (audit~A22).

\section{What remains}\label{sec:open}

\begin{enumerate}
\item \emph{$\PTM_2$ beyond girth five.} The walk theorem reduces it to the base separation
problem; Theorem~\ref{thm:sparse} solves that problem only under sparseness. The obstructions to
sparseness are short cycles through the special vertices. $K_4$ shows that some such
configurations are genuinely bad.
\item \emph{More than two edges.} The walk theorem chains: leaf the first edge, and for each
further edge walk the even vertex from the previous short end to the next long end, deleting
the walk; restore in reverse order, re-attaching each leaf before walking back. The base graph
then needs separations of $\{x_i,y_i\}$, $\{x_1,\ell_1\}$ and all pairs of long ends, and each
walk must avoid the neighbourhoods of the later short ends. Writing down a clean sufficient
condition (large girth and large mutual distances) is routine but has not been done or
formalised.
\item \emph{The multigraph case.} Theorem~\ref{thm:walk} assumes $G$ simple. For STL at a vertex
whose neighbourhood is dense, every pairing may contain adjacent pairs.
\item \emph{Relation to EE.} $\PTM_2(r)$ at $d(z)=4$ contains the end-edge conjecture EE of~\cite{PE}
at those $u\in N(z)$ that lie in a pairing $\{uu',cd\}$ of $N(z)$ with $uu',cd\notin E(J)$: the
decomposition given by $\PTM(G,M;uu'@u)$ lifts to one in which $uz$ is terminal at $u$. For the
other $u$ the graph $(J-z)+M$ is a multigraph, outside Definition~\ref{def:ptm}. So a full proof
of $\PTM_2(r)$ including the multigraph case would settle the open case of EE
($|N(u)\cap N(z)|\ge2$) at $d(z)=4$. The walk theorem's final
step is EE-like (the Fan transfers end at the long end $y$); the girth hypothesis plays the role
of ``at most one common neighbour'' in~\cite[Theorem~5.1]{PE}.
\end{enumerate}

\section{The formalisation}\label{sec:lean}

\texttt{lean/GallaiLean/TerminalMatching.lean} (imports \texttt{StarTerminal}) contains, without
\texttt{sorry}; \texttt{manuscripts/terminal-matching/PTMAxioms.lean} prints for the main
theorems only \texttt{propext}, \texttt{Classical.choice}, \texttt{Quot.sound}, and the full
library builds.
\begin{itemize}
\item \texttt{chainEdges}, \texttt{ChainAdj}, \texttt{delChain}, \texttt{delChain\_degree\_mod}:
paths as vertex lists, and the parity of degrees after deleting the edges of a path.
\item \texttt{walk\_back}: Lemma~\ref{lem:walkback}, by induction along the list, each step
\texttt{fan\_single\_supp}.
\item \texttt{leafG}, \texttt{baseG}, \texttt{baseG\_parity}, \texttt{ptm2\_of\_base}:
Definition~\ref{def:walk} and Theorem~\ref{thm:walk}; the conclusion is \texttt{PTM2At G a b x y}.
The two leaves live in \texttt{Option (Option W)}; \texttt{PathDecomp.liftLeaf} and
\texttt{leaf\_merge} (twice) do the leaf bookkeeping.
\item \texttt{sep\_sparse\_fix}, \texttt{ptm2\_of\_sparse}: Lemma~\ref{lem:fix} and
Theorem~\ref{thm:sparse}, from \texttt{separation\_sparse}.
\item \texttt{NoShortCycleAt}, \texttt{ptm2\_girth\_core}, \texttt{ptm2\_localgirth}, and
\texttt{ptm2\_girth5} through \texttt{NoTriangle}, \texttt{NoSquare}: Theorem~\ref{thm:girth}. The minimal walk is taken by \texttt{Nat.find} and made a path by
\texttt{Walk.bypass}; the distance arguments use \texttt{takeUntil}/\texttt{dropUntil}.
\item \texttt{swapG}, \texttt{liftJ}, \texttt{swapG\_degree}, \texttt{GPath.subdivide\_start},
\texttt{PathDecomp.replace\_edges}, \texttt{unlift}, \texttt{stl\_of\_ptm2}: Theorem~\ref{thm:stl}.
\item \texttt{stl\_of\_localgirth}, \texttt{test1\_of\_localgirth} (and the girth-five forms
\texttt{stl\_of\_girth5}, \texttt{test1\_of\_girth5}): Corollary~\ref{cor:test1}, through
\texttt{test1\_of\_stl}.
\end{itemize}

\begin{thebibliography}{9}
\bibitem{CONN} Project notes, \texttt{research/connections.md}, \S\S X, Y, Z.
\bibitem{Fan05} G.~Fan, Path decompositions and Gallai's conjecture, J.~Combin.\ Theory Ser.~B 93
(2005) 117--125.
\bibitem{L68} L.~Lov\'asz, On covering of graphs, in: Theory of Graphs (Proc.\ Colloq.\ Tihany,
1966), Academic Press, 1968, 231--236.
\bibitem{MS} J.~Williams, Separating many pairs in Lov\'asz decompositions: classes, the leaf
merge, and the Case~I lemma, \texttt{manuscripts/multi-sep}, 2026.
\bibitem{PE} J.~Williams, Prescribed end edges in Lov\'asz decompositions with one even vertex,
\texttt{manuscripts/prescribed-end}, 2026.
\bibitem{RD} J.~Williams, Route descent, \texttt{manuscripts/route-descent}, 2026.
\bibitem{SP} J.~Williams, The spine: Gallai's conjecture and two deletion tests on $3$-connected
graphs, \texttt{manuscripts/spine}, 2026.
\bibitem{STL} J.~Williams, Star-terminal Lov\'asz decompositions and the deletion tests for
Gallai's conjecture, \texttt{manuscripts/stl}, 2026.
\end{thebibliography}
\end{document}
